\section{Coverage before cancellation in the scalar self-seeded basis}
\label{sec:coverage-no-cancellation}

This section isolates a mathematical simplification of the supplied
scalar self-seeded presentation. It is not a complete new basis below
the measured checker frontier. In particular, it does not justify
deleting successor injectivity from that presentation.

Use ordinary intuitionistic first-order logic with equality, terms
$0,S,\epsilon$, and a staged scalar recursor $F(n;v.h)$ with equations
\[
 F(0;v.h)=S(h[0/v]),\qquad
 F(Sn;v.h)=h[F(n;v.h)/v].
\]
The remaining nonlogical assumptions, besides contextual induction on
literal equations, are explicitly
\[
 0=S0\to\bot,\qquad
 \epsilon(a,b)=0\lor\epsilon(a,b)=1,\qquad
 \epsilon(a,b)=1\leftrightarrow a=b.
\]
All free numerical variables in these schemas are universally closed. Here $1=S0$, and each recursor body is an earlier-stage
source term. Neither successor injectivity nor numerical coverage is
assumed in the result below.

Put $D(c,a)=F(c;v.a)$ and $Z(c)=D(c,0)$. Thus
\[
 D(0,a)=Sa,\quad D(Sn,a)=a,\quad Z(0)=1,\quad Z(Sn)=0.
\]
If $Sn=0$, these equations give $0=1$, so every successor is nonzero.
For Boolean $b,c$, the source terms
\[
 \neg_Bb=\epsilon(b,0),\qquad
 b\wedge_Bc=\epsilon(D(b,\epsilon(c,0)),0)
\]
have their usual truth tables; define $\vee_B$ by De Morgan. Only
$0\ne1$ and successor nonzeroness are used in these finite cases.
The usual external syntax recursion assigns a Boolean characteristic
to each fixed quantifier-free formula. Literal-equation induction
therefore gives contextual quantifier-free induction, without
successor cancellation.

Define the following fixed source terms:
\[
 \begin{split}
 B_a(v)&=\epsilon(v,0)\vee_B\epsilon(Sv,a),\\
 H_a(v)&=D(B_a(v),v),\\
 U(j,a)&=F(j;v.H_a(v)),\qquad p_*(a)=U(a,a).
 \end{split}
\]
The update holds the state at zero or immediately before $a$, and
increments it in the remaining case. Since $B_a(0)=1$,
$H_a(0)=0$ and $U(0,a)=1$. With $a$ fixed, quantifier-free induction gives
\begin{equation}
 U(j,a)=Sj\quad\lor\quad S(U(j,a))=a.
 \label{eq:protected-orbit}
\end{equation}
At zero the first alternative holds. A state satisfying the second
alternative is fixed by $H_a$. Otherwise the first alternative advances
from $Sj$ to $S(Sj)$: in that branch the current state is a successor
and hence nonzero. No separate induction asserting nonzeroness of the
whole orbit is needed.

After this induction has been proved with its parameter fixed,
specialize $j=m$ and $a=S^2m$. Both alternatives of
\eqref{eq:protected-orbit} imply $S(U(m,S^2m))=S^2m$. The state is
then fixed. Exactly two further computation steps give
\begin{equation}
                     S(p_*(S^2m))=S^2m.
 \label{eq:protected-right-inverse}
\end{equation}
This is a right-inverse equation only; no cancellation of $S$ has
taken place.

\paragraph{Whole-domain coverage theorem.}
The theory with the assumptions just listed proves
\[
                     n=0\ \lor\ \exists m\,n=Sm.
\]
Indeed, apply quantifier-free induction to
\[
 C(n):\quad n=0\ \lor\ n=1\ \lor\ S(p_*(n))=n.
\]
The zero case is immediate. In the successor step, $n=0$ gives
$Sn=1$. If $n=1$, apply \eqref{eq:protected-right-inverse} with
$m=0$. In the third case $n=S(p_*(n))$, whence
$Sn=S^2(p_*(n))$; apply \eqref{eq:protected-right-inverse} with
$m=p_*(n)$ and rewrite. This proves $C(n)$ for every $n$.
Its second and third alternatives provide the witnesses $0$ and
$p_*(n)$ respectively. The proof covers nonstandard elements and
never assumes that an arbitrary element is reached by a standard
finite successor iteration.

\subsection{No fixed points and rigidity of each standard numeral}

The same weak theory proves more than coverage, although the following
results still fall short of arbitrary successor cancellation.

\paragraph{Successor has no fixed point.}
Let $h(v)=\epsilon(v,0)$. This source update has no fixed point.
Indeed, $h(v)=v$ and the Boolean range of $h$ give $v=0$ or $v=1$;
the two cases contradict $h(0)=1$ and $h(1)=0$, respectively.
If $Sa=a$, the recursor equation would give
\[
 F(a;v.h)=F(Sa;v.h)=h(F(a;v.h)),
\]
contrary to that fact. Thus $Sa\ne a$ for every $a$, without using
successor injectivity or coverage. In particular, $1\ne2$.

\paragraph{The only predecessor of one is zero.}
Put
\[
 h(v)=\neg_B\epsilon(v,2),\qquad A(n)=F(n;v.h).
\]
The distinctness of $0,1,2$ gives $h(0)=1$, $A(0)=2$, and
$A(1)=h(2)=0$. The image of $h$ is contained in $\{0,1\}$, and
\[
                         h(v)=0\quad\longleftrightarrow\quad v=2.
\]
Quantifier-free induction proves $A(n)=2\leftrightarrow n=0$.
The zero case holds, whereas at every successor both sides are false:
$A(Sn)=h(A(n))$ is zero or one, and $Sn\ne0$.
If $Sn=1$, then $h(A(n))=A(Sn)=A(1)=0$; hence $A(n)=2$ and $n=0$.
The converse is immediate. Therefore
\[
                              Sn=1\quad\longleftrightarrow\quad n=0.
\]

\paragraph{Finite-prefix rigidity metatheorem.}
For every fixed standard integer $k\ge1$, there is a finite derivation
in the weak theory of
\begin{equation}
 Sx=\overline{k}\quad\longleftrightarrow\quad x=\overline{k-1}.
 \label{eq:standard-predecessor}
\end{equation}
Here bars denote fixed numerals. This is an external metatheorem,
not a derivation of a single formula quantifying over $k$.

The proof is an external induction on $k$. Carry as induction data
the distinctness of the numerals $0,\ldots,k+1$ and
\eqref{eq:standard-predecessor} for all smaller positive indices.
The initial distinctness of $0,1,2$ follows above. Define the following
finite numeral-valued source update:
\[
 h_k(v)=
 \begin{cases}
 \overline{i-1},&v=\overline{i},\quad 1\le i\le k-1,\\
 \overline{k-1},&v=\overline{k+1},\\
 \overline{k},&\text{otherwise}.
 \end{cases}
\]
This display abbreviates an actual source term, not an additional
recursion schema over a graph-defined function. To construct it,
for each fixed $r=1,\ldots,k$ form the Boolean characteristic $b_r(v)$
of the finite condition that the displayed value is at least $r$.
Explicitly, $b_r$ is the Boolean negation of the disjunction of
$\epsilon(v,\overline{i})$ for $1\le i\le\min(r,k-1)$, together
with $\epsilon(v,\overline{k+1})$ when $r=k$. Empty disjunctions
mean zero; all index comparisons here are metalevel comparisons
between fixed integers, not assumed arithmetic order predicates.
Start with $t_0(v)=0$ and set
\[
                  t_r(v)=D(\neg_B b_r(v),t_{r-1}(v)).
\]
Then $h_k=t_k$: each true condition contributes one successor, and
each false condition leaves the value unchanged. The construction
uses exactly $k$ external syntax steps and the already available
finite Boolean gates. It is not asserted to provide a single source
term with variable $k$.

Write $A_k(n)=F(n;v.h_k)$ and introduce the distinct marker numerals
\[
 q_0=\overline{k+1},\qquad q_j=\overline{k-j}\quad(1\le j\le k).
\]
Since $h_k(0)=\overline{k}$, the initial state is $q_0$. Direct finite
computation gives $A_k(\overline{j})=q_j$ for $0\le j\le k$.
The displayed update has no output $q_0$, and for $1\le j\le k$ it
satisfies
\[
                       h_k(v)=q_j\quad\longleftrightarrow\quad v=q_{j-1}.
\]
Quantifier-free induction therefore gives
$A_k(n)=q_0\leftrightarrow n=0$. For $1\le j<k$, an external
induction on $j$, using numerical coverage and the already proved
predecessor uniqueness for $\overline{j}$, gives
\[
                     A_k(n)=q_j\quad\longleftrightarrow\quad n=\overline{j}.
\]
For clarity, in the successor case $n=Sm$ the left side is equivalent
to $A_k(m)=q_{j-1}$, hence to $m=\overline{j-1}$. The right side has
exactly the same condition by the previous predecessor theorem.
The zero case has two false sides.

Now $Sx=\overline{k}$ implies
$h_k(A_k(x))=A_k(\overline{k})=q_k$, so $A_k(x)=q_{k-1}$ and
$x=\overline{k-1}$. This proves the next predecessor theorem.
To extend numeral distinctness to $k+2$, a putative equality
$\overline{k+2}=\overline{j}$ with $1\le j\le k$ reduces, by the
newly available predecessor theorems, to
$\overline{k+1}=\overline{j-1}$, already impossible. Equality with
zero is excluded by successor nonzeroness, and equality with $k+1$
by the no-fixed-point lemma. This closes the external induction.

Consequently no element other than the intended previous numeral can
map to any fixed standard positive numeral. This does not rule out
collisions over nonstandard targets: replacing the fixed construction
above by a uniform variable-index argument would require an additional
theorem that has not been supplied.

\paragraph{A simpler predecessor when cancellation is retained.}
If successor injectivity is retained, as in the supplied complete
self-seeded theory, \eqref{eq:protected-right-inverse} gives
$p_*(S^2m)=Sm$. Define the total single-valued finite-case graph
\[
 P(n)=
 \begin{cases}
 0,&n=0\text{ or }n=1,\\
 p_*(n),&n\ne0\text{ and }n\ne1.
 \end{cases}
\]
Its graph uses only quantifier-free source guards and source value
terms. Finite common refinements therefore eliminate each fixed
quantifier-free formula containing $P$, preserving the available
induction rule. The zero equation is immediate. For $P(Sn)=n$,
use the already derived coverage to split $n=0$ or $n=Sm$.
The first gives $P1=0$. In the second, $S^2m$ is neither zero nor
one, and the displayed equation gives $P(S^2m)=Sm=n$.

Thus only the values at zero and one need correction, rather than
the three exceptional values in the supplied stopped-counter
construction. Coverage is proved first, and the subsequent
predecessor proof no longer needs a joint coverage/predecessor
induction. The remaining open issue is precise: can the weaker
theory prove $p_*(S^2m)=Sm$, or otherwise prove $Sx=Sy\to x=y$?
The proof above establishes neither statement.
